正式矩阵包含120个任务，其中28/30个订单流的四种机制均可执行；完整机制在15/30条可执行订单流中实际发生跨场服务，并在10/28个完整配对中满足全日参与底线。其中2个受控不可执行单元为50客户—空间交错—流1—无参与底线、50客户—空间交错—流3—无参与底线；这些单元保留在可执行率中，不进入四种机制的成对均值。完整机制有2个阶段超过下一触发间隔，可比较阶段共54个，最长阶段用时36789.5 s；因此本批结果只能解释为批量滚动决策支持，不能改称实时求解。

在28个完整配对中，完整机制相对禁合作的净收益差均值为+58.9，改善/变差/持平为16/12/0；收入、成本、完成客户数和完成需求差均值分别为+2.3、-56.6、+0.04和+12.3。最有利单元是150客户—空间交错—流4（+634.6），最不利单元是150客户—地理聚集—流4（-351.1）。净收益差必须与服务收入和工作量差共同解释；这里报告的是观察到的全日结果，不把均值方向直接解释成单一机制的因果效应。

在28个完整配对中，完整机制相对无参与底线的净收益差均值为-26.3，改善/变差/持平为6/14/8；收入、成本、完成客户数和完成需求差均值分别为+0.2、+26.6、+0.00和+1.2。最有利单元是150客户—空间交错—流2（+577.2），最不利单元是150客户—地理聚集—流4（-529.4）。净收益差必须与服务收入和工作量差共同解释；这里报告的是观察到的全日结果，不把均值方向直接解释成单一机制的因果效应。

在28个完整配对中，完整机制相对顺序插入基线的净收益差均值为+193.2，改善/变差/持平为25/3/0；收入、成本、完成客户数和完成需求差均值分别为-9.6、-202.8、-0.07和-50.7。最有利单元是100客户—地理聚集—流5（+614.7），最不利单元是150客户—地理聚集—流4（-323.4）。净收益差必须与服务收入和工作量差共同解释；这里报告的是观察到的全日结果，不把均值方向直接解释成单一机制的因果效应。

28个电网日的充电择时复算形成840组配对，改善/变差/持平为405/364/71。六个网络—责任单元的充电排放降幅介于2.78\%和3.28\%之间，最低出现在50客户—空间交错，最高出现在100客户—地理聚集；总运营排放降幅范围为0.29\%--0.39\%。该结果固定路径、车辆、服务客户和充电电量，只识别合法窗口内充电择时的增量作用，不构成路径—充电联合优化的证据。
